\section{Appendix on incidence ingredients for Sections \ref{sec:subsidy} and \ref{sec:decentralized}}\label{app:subsidy}

\subsection{Mapping from product-market primitives to local welfare}

This appendix collects several useful decompositions for the local incidence block used in Section \ref{sec:decentralized}, building on the primitive benchmark introduced in Section \ref{sec:subsidy}.

Recall that local value under independent ownership is
\begin{equation}
    B_j^{L,D}
    =
    C^D(j)
    + \lambda_I^L \pi_I^D(j)
    + \lambda_E^L \pi_E^D(j)
    + \eta_j^D,
\end{equation}
whereas local value under approved acquisition is
\begin{equation}
    B_j^{L,A}
    =
    C^M(j)
    + \lambda_M^L \Pi^M(j)
    + \eta_j^A.
\end{equation}
Hence the local effect of acquisition relative to independent ownership is
\begin{equation}\label{eq:localdifference}
    B_j^{L,A}-B_j^{L,D}
    =
    \bigl(C^M(j)-C^D(j)\bigr)
    +
    \bigl(\lambda_M^L \Pi^M(j)-\lambda_I^L \pi_I^D(j)-\lambda_E^L \pi_E^D(j)\bigr)
    +
    \bigl(\eta_j^A-\eta_j^D\bigr).
\end{equation}
Equation \eqref{eq:localdifference} makes clear that approved acquisition can reduce local value even if it raises private surplus.
A sufficient condition for \(B_j^{L,A}<B_j^{L,D}\) is that the local loss in consumer surplus and spillovers dominates the locally retained component of merged-firm profits.

\subsection{A decomposition with locally retained buyout proceeds}\label{app:localbuyoutproceeds}

The main text emphasizes that local governments may value successful buyout-oriented entrepreneurship partly because some gains from exit are retained within the locality.
This appendix makes that interpretation explicit.

Suppose that, under approved acquisition, local value can be decomposed as
\begin{equation}\label{eq:BLAtransfer}
    B_j^{L,A}
    =
    C^M(j)
    + \lambda_M^L \Pi^M(j)
    + \mu_j^L T_j
    + \eta_j^A,
\end{equation}
where \(T_j \ge 0\) denotes the buyout payment associated with a successful type-\(j\) acquisition and
\(\mu_j^L \in [0,1]\) is the fraction of that payment that is retained locally.
For example, \(\mu_j^L\) may reflect the local residence of founders, employees with equity claims, or locally taxed capital gains.

Under independent ownership, local value remains
\begin{equation}\label{eq:BLDtransfer}
    B_j^{L,D}
    =
    C^D(j)
    + \lambda_I^L \pi_I^D(j)
    + \lambda_E^L \pi_E^D(j)
    + \eta_j^D.
\end{equation}

By contrast, from the national perspective the buyout payment \(T_j\) is a transfer rather than a net surplus component.
A corresponding national decomposition is therefore
\begin{equation}\label{eq:BNAtransfer}
    B_j^{N,A}
    =
    C^M(j)
    + \Pi^M(j)
    + \xi_j^A,
\end{equation}
with no separate \(T_j\) term.

Under this interpretation, the local--national difference under approved acquisition is
\begin{equation}\label{eq:approvedwedgewithtransfer}
    B_j^{L,A} - B_j^{N,A}
    =
    (\lambda_M^L-1)\Pi^M(j)
    + \mu_j^L T_j
    + \eta_j^A - \xi_j^A.
\end{equation}
Equation \eqref{eq:approvedwedgewithtransfer} clarifies why buyout-oriented projects may be valued differently by local and national planners.
Locally retained buyout proceeds, represented by \(\mu_j^L T_j\), can raise local value even when they do not increase national surplus.

Combining \eqref{eq:BLAtransfer} and \eqref{eq:BLDtransfer}, the local effect of approved acquisition relative to independent ownership becomes
\begin{equation}\label{eq:localdifferencewithtransfer}
    B_j^{L,A}-B_j^{L,D}
    =
    \bigl(C^M(j)-C^D(j)\bigr)
    + \bigl(\lambda_M^L\Pi^M(j)-\lambda_I^L\pi_I^D(j)-\lambda_E^L\pi_E^D(j)\bigr)
    + \mu_j^L T_j
    + \bigl(\eta_j^A-\eta_j^D\bigr).
\end{equation}
Hence approved acquisition can be locally attractive even when it reduces local consumer surplus, provided the locally retained buyout proceeds and other locally retained gains are sufficiently large.

This decomposition also helps interpret the equilibrium incidence wedge in Section \ref{sec:decentralized}.
If buyout-oriented projects have relatively large \(T_j\), then the local government may value successful acquisition more highly than the national planner, even though the underlying product-market outcome is the same.

\subsection{Expected local value under endogenous merger attempt}

Let
\begin{equation}
    a_j(q)=\mathbf{1}\{q\Delta_j \ge k\}
\end{equation}
be the merger-attempt indicator.
If \(a_j(q)=0\), then successful entry necessarily yields independent ownership, so
\begin{equation}
    B_j^L(q)=B_j^{L,D}.
\end{equation}
If \(a_j(q)=1\), then the merger is attempted and successful entry yields approved acquisition with probability \(q\) and independent competition with probability \(1-q\).
The expected local value is therefore
\begin{equation}
    q B_j^{L,A} + (1-q)B_j^{L,D} - \kappa^L k
    =
    B_j^{L,D}
    +
    q\bigl(B_j^{L,A}-B_j^{L,D}\bigr)
    -
    \kappa^L k.
\end{equation}
Combining the two cases gives
\begin{equation}
    B_j^L(q)
    =
    B_j^{L,D}
    +
    a_j(q)
    \left[
        q\bigl(B_j^{L,A}-B_j^{L,D}\bigr)
        -
        \kappa^L k
    \right].
\end{equation}

\subsection{Comparative statics of local value}

On any interval where \(q<\bar q_j \equiv k/\Delta_j\), the merger is not attempted, so
\begin{equation}
    \frac{d B_j^L(q)}{dq}=0.
\end{equation}
On any interval where \(q>\bar q_j\), the merger is attempted, so
\begin{equation}
    \frac{d B_j^L(q)}{dq}
    =
    B_j^{L,A}-B_j^{L,D}.
\end{equation}
Hence if approved acquisition lowers local value relative to independent ownership, then \(B_j^L(q)\) is weakly decreasing in \(q\) once merger attempt becomes privately viable.

\subsection{A sufficient condition for subsidy withdrawal after project switching}

Suppose the buyout-oriented project becomes privately optimal at \(q=\hat q\).
Because the local planner solves the same quadratic subsidy problem with \(B_B^{PM}(q)\) replaced by \(B_B^L(q)\), the local subsidy falls to zero immediately to the right of \(\hat q\) if
\begin{equation}
    B_B^L(\hat q)
    <
    \frac{\tau+2}{2}R_B(\hat q).
\end{equation}
Using the definition of \(B_B^L(\hat q)\), this condition can be written as
\begin{equation}
    B_B^{L,D}
    +
    \hat q\bigl(B_B^{L,A}-B_B^{L,D}\bigr)
    -
    \kappa^L k
    <
    \frac{\tau+2}{2}
    \left[
        \pi_E^D(B)+\beta(\hat q\Delta_B-k)
    \right],
\end{equation}
because project \(B\) is merger-feasible at \(q=\hat q\) under Assumption \ref{ass:singlecross}.
This decomposition is useful for interpretation:
subsidy withdrawal is more likely when approved acquisition strongly reduces local value, when attempt costs are large, or when the entrepreneur's private buyout return is especially high.

\subsection{General distribution of entry costs}

This appendix extends the local-subsidy problem to a general entry-cost distribution.

Let fixed costs be drawn from a continuously differentiable CDF \(G\) with density \(g\).
For a given project type \(j\), the entry cutoff remains
\begin{equation}
    c_j(s,q)=s+\frac{R_j(q)^2}{2}.
\end{equation}
Aggregate local welfare is
\begin{equation}\label{eq:WgeneralGlocal}
    W_{j,G}^L(s,q)
    =
    \int_0^{c_j(s,q)}
    \left[
        R_j(q)B_j^L(q)
        -
        \frac{R_j(q)^2}{2}
        -
        f
        -
        \tau s
    \right] dG(f).
\end{equation}
Equivalently,
\begin{equation}
    W_{j,G}^L(s,q)
    =
    \left(
        R_j(q)B_j^L(q)
        -
        \frac{R_j(q)^2}{2}
        -
        \tau s
    \right)
    G(c_j)
    -
    \int_0^{c_j} f\,dG(f).
\end{equation}
Differentiating with respect to \(s\), using \(dc_j/ds=1\), yields the first-order condition
\begin{equation}\label{eq:focgeneralGlocal}
    g(c_j)
    \left[
        R_j(q)B_j^L(q)
        -
        R_j(q)^2
        -
        (\tau+1)s
    \right]
    -
    \tau G(c_j)
    =
    0.
\end{equation}
The uniform-distribution expressions used in Sections \ref{sec:subsidy} and \ref{sec:decentralized} are recovered as the special case \(G(f)=f\).

\subsection{Why the revised welfare block differs from the earlier reduced-form one}

The earlier reduced-form formulation treated the successful-project value as a direct convex combination of independent ownership and approved acquisition.
That shortcut obscures an important equilibrium fact:
if \(q\Delta_j<k\), no merger is attempted, so a change in \(q\) cannot affect realized outcomes for project \(j\).

The revised welfare object corrects this by inserting the attempt indicator \(a_j(q)\).
As a result, merger permissiveness affects local subsidy design only after it changes actual behavior.
This is the sense in which the revised welfare block is behaviorally consistent with the private equilibrium derived in Section \ref{sec:private}.
